% CS-433 Machine Learning, Fall 2026 - Project 1 report
% Adapted from the course template:
% https://github.com/epfml/ML_course/tree/main/projects/project1/latex-example-paper
%
% Hard limit: 2 pages of content + 1 extra page for references only.
% Compile with:
%   pdflatex report && bibtex report && pdflatex report && pdflatex report

\documentclass[10pt,conference,compsocconf]{IEEEtran}

\usepackage{hyperref}
\usepackage{graphicx}  % for \includegraphics
\usepackage{amsmath}   % for math environments

\begin{document}

\title{Predicting Coronary Heart Disease Risk from BRFSS Lifestyle Data}

\author{
  Name1 Surname1, Name2 Surname2, Name3 Surname3\\
  \textit{Department of Computer Science, EPFL, Switzerland}
}

\maketitle

\begin{abstract}
  % 4 things to cover: the problem, why it matters, what your solution
  % achieves, what follows from it. Write this section last.
\end{abstract}

\section{Introduction}
\label{sec:introduction}
% State the problem (predicting MICHD from BRFSS survey data), why it is
% interesting/relevant, and summarize your contributions/approach.

\section{Data and Exploratory Analysis}
\label{sec:eda}
% Describe the dataset (size, features, class imbalance), what you found in
% exploratory data analysis, and how it motivated your preprocessing choices.
% See experiments/eda.py.

\section{Models and Methods}
\label{sec:methods}
% Describe your preprocessing/feature engineering pipeline (src/preprocessing.py,
% src/features.py) and the models you used, including the 6 required methods
% from implementations.py and any custom extension (src/models.py). Include
% enough detail (hyperparameters, feature transforms, number of CV folds, ...)
% for reproducibility.

% Example figure:
% \begin{figure}[tbp]
%   \centering
%   \includegraphics[width=\columnwidth]{figures/my_figure}
%   \caption{Caption describing the figure.}
%   \label{fig:my-figure}
% \end{figure}

\section{Results}
\label{sec:results}
% Report cross-validated performance of the baseline methods (Table 1: the 6
% required methods) and of your improved pipeline. Include the ablation study:
% which modification helped the most? See experiments/baselines.py,
% experiments/tuning.py, experiments/ablation.py.

% Example table:
% \begin{table}[htbp]
%   \centering
%   \begin{tabular}{|l|c|c|}
%     \hline
%     Method & CV Accuracy & CV F1 \\
%     \hline
%     Least squares & & \\
%     Ridge regression & & \\
%     Logistic regression & & \\
%     Our final model & & \\
%     \hline
%   \end{tabular}
%   \caption{Cross-validated performance comparison.}
%   \label{tab:results}
% \end{table}

\section{Discussion}
\label{sec:discussion}
% Discuss strengths/weaknesses of your approach based on the results above.
% Explain unexpected results, failure cases, limitations.

\section{Summary}
\label{sec:summary}
% Summarize your contributions and findings.

\bibliographystyle{IEEEtran}
\bibliography{references}

\end{document}
